\documentclass[10pt,letterpaper]{article}
\usepackage{cornell-notes}

\title{Clause 33.06.04.03.02: Class Timed Mutex}
\author{}
\date{}
\setDocTitle{Clause 33.06.04.03.02: Class Timed Mutex}
\setDocSubtitle{C++ 2024}
\setDocOwner{C++ 2024 Cornell Notes}
\setDocDate{\today}
\setCornellCollection{C++ 2024 Cornell Notes}
\setCornellUnitType{Clause}
\setCornellUnitNumber{33.06.04.03.02}
\setCornellUnitTitle{Class Timed Mutex}
\hypersetup{pdftitle={Clause 33.06.04.03.02: Class Timed Mutex},pdfsubject={Cornell notes},pdfauthor={}}

\begin{document}
\maketitle
\begin{CornellOverviewBox}{Overview}
\textbf{Classification:} Standard library - Normative\\[2pt]
\textbf{Learning objective:} Understand the rules, terminology, and semantic consequences associated with Class timed\_mutex.
\end{CornellOverviewBox}\begin{CornellSummaryBox}{Summary}
{\sffamily\bfseries\color{Accent} Summary}\par\vspace{3pt}Section 33.6.4.3.2 focuses on Class timed\_mutex. Use the listed grammar productions together with the semantic constraints and disambiguation rules referenced by the section. When applying it, preserve the distinction between mandatory requirements, recommendations, permissions, and behavior deliberately left undefined, unspecified, or implementation-defined.
\end{CornellSummaryBox}

\end{document}
